\chapter{Gaussian quadratic forms}\label{ch:gaussian}
\chapterlead{A quadratic statistic may be presented as a matrix, a factor, or a list of
spectral coefficients. The representation determines how the calculation is
organized. Under the Gaussian model, the spectrum determines the law.
This chapter fixes that distinction before any numerical approximation is
introduced.}
\label{sec:gaussian-core}

The primary mathematical object is a Gaussian quadratic form.

\section{Central positive forms}

Let
\[
X\sim N(0,\Sigma),
\qquad
M\succeq0,
\]
and write \(S=\Sigma^{1/2}\).  With \(X=SZ\), \(Z\sim N(0,I)\),
\[
Q=X^\top MX
=
Z^\top(SMS)Z.
\]
If the nonzero eigenvalues of the positive semidefinite operator
\[
H=SMS
\]
are \(\lambda_1,\ldots,\lambda_r\), then
\[
\boxed{
Q\overset d=
\sum_{j=1}^r\lambda_jZ_j^2.
}
\]

Here and throughout the central reductions, the variables $Z_j$ are independent
standard normal variables. The numerical tail problem therefore depends only
on a positive spectrum, regardless of how that spectrum was generated.
If there are no nonzero coefficients, then $Q=0$ almost surely and
$P(Q\ge q)=1$ for $q\le0$, while $P(Q\ge q)=0$ for $q>0$.
This deterministic case is handled directly, without a saddlepoint or an
inversion calculation.

\begin{proposition}[Representation equivalence]\label{claim:final-2-1}
The same central quadratic-form law may be supplied by any of the following:
\begin{enumerate}
\item a nonnegative spectrum \(\{\lambda_j\}\).
\item a positive semidefinite matrix \(H\).
\item a factor \(A\) with \(H=AA^\top\) or \(H=A^\top A\).
\item a replayable streamed factor
\[
A=[A_1\ \cdots\ A_B].
\]
\end{enumerate}
The nonzero spectra of \(AA^\top\) and \(A^\top A\) are identical, and
\[
AA^\top=\sum_bA_bA_b^\top.
\]
\end{proposition}

\begin{proof}
If $AA^\top u=\lambda u$ with $\lambda>0$, then $A^\top u\ne0$ and
$A^\top A(A^\top u)=\lambda A^\top u$. Conversely, an eigenvector
$v$ of $A^\top A$ with positive eigenvalue maps to the eigenvector $Av$
of $AA^\top$. These maps restrict to inverse isomorphisms of the
corresponding eigenspaces after scaling, so multiplicities agree.
Horizontal concatenation gives $AA^\top=\sum_b A_bA_b^\top$ by direct
matrix multiplication. Orthogonal diagonalization then identifies the
associated Gaussian quadratic-form laws.
\end{proof}

This proposition is the domain-independent source of the memory-bounded
streaming method developed below.

\section{Central indefinite forms}

The positive case is not the natural boundary of the central theory.  Let
\[
X\sim N(0,\Sigma),
\qquad
M=M^\top,
\]
where $M$ may be indefinite, and again write $S=\Sigma^{1/2}$.  The symmetric
operator
\[
H=SMS
\]
has a real signed spectrum.

\begin{theorem}[Exact central signed reduction]\label{claim:final-2-2}
If $\lambda_1,\ldots,\lambda_r$ are the nonzero eigenvalues of $H$, then
\[
\boxed{
X^\top M X
\overset d=
\sum_{j=1}^r\lambda_jZ_j^2.
}
\]
\end{theorem}

\begin{proof}
Write $X=SZ$ with $Z\sim N(0,I)$ and diagonalize
\[
SMS=V\Lambda V^\top.
\]
Rotational invariance of the standard normal gives
\[
Z^\top V\Lambda V^\top Z
\overset d=
\sum_j\lambda_jZ_j^2.
\]
No positivity assumption on $M$ is needed for this central reduction.
\end{proof}

The sign pattern determines the support.  If every $\lambda_j\ge0$, the lower
endpoint is zero.  If every $\lambda_j\le0$, the upper endpoint is zero.  If
both signs occur, the distribution is unbounded in both directions.

For the right-tail path, the cumulant generating function is
\[
\kappa(s)
=
-\frac12\sum_j\log(1-2s\lambda_j),
\]
with real domain
\[
\boxed{
\max_{\lambda_j<0}\frac{1}{2\lambda_j}
< s <
\min_{\lambda_j>0}\frac{1}{2\lambda_j}.
}
\]
Here a missing negative or positive spectral set gives a domain endpoint
\(-\infty\) or \(+\infty\), respectively.
When $q>E[Q]=\sum_j\lambda_j$, the saddlepoint is the unique positive solution
of
\[
\kappa'(\theta)=q
\]
inside this domain whenever a positive spectral component is present.  The
negative-only case with at least one nonzero coefficient has a positive
saddlepoint only for $E[Q]<q<0$. Indeed, $\kappa'(\theta)$ increases from
$E[Q]$ to zero as $\theta\to\infty$. For $q\ge0$ its survival
probability is zero, and no finite positive saddlepoint solves the query.

The packet representation is sign aware: positive and negative eigenvalues are
clustered separately, and each same-sign packet stores a center $\mu$ and
centered moments.  The local expansion remains
\[
-\frac12\log(1-2s(\mu+\delta))
=
-\frac12\log(1-2s\mu)
+
\frac12\sum_{n\ge1}
\frac{\alpha(s)^n\delta^n}{n},
\qquad
\alpha(s)=\frac{2s}{1-2s\mu},
\]
with the same geometric remainder condition
\[
|\alpha(s)|\max|\delta|<1.
\]

A useful distinction appears in the Fourier tail.  Under the saddlepoint tilt,
\[
\left|\phi_\theta(t)\right|
=
\prod_j
\left[
1+
\left(
\frac{2\lambda_j t}{1-2\theta\lambda_j}
\right)^2
\right]^{-1/4}.
\]
If $r$ coefficients are nonzero, then asymptotically
\[
|\phi_\theta(t)|=O(|t|^{-r/2}).
\]
Thus certification difficulty depends strongly on spectral rank, not only on
how small the requested probability is.

The signed law has an exact spectral reduction, but its numerical
certification requires bounds appropriate to its support and Fourier decay.
The finite-endpoint, mixed-sign, and low-rank controls are developed in
Chapter~\ref{ch:evidence}.

\section{Noncentral positive forms}

Let
\[
X\sim N(\mu,\Sigma),
\qquad
M\succeq0.
\]
Again set \(S=\Sigma^{1/2}\) and diagonalize
\[
SMS=V\Lambda V^\top.
\]
Writing
\[
c=V^\top SM\mu,
\]
each positive-eigenvalue direction satisfies
\[
\lambda_iY_i^2+2c_iY_i
=
\lambda_i
\left(
Y_i+\frac{c_i}{\lambda_i}
\right)^2
-
\frac{c_i^2}{\lambda_i}.
\]

\begin{theorem}[Exact noncentral reduction]\label{claim:final-2-3}
For the positive eigenvalues \(\lambda_i>0\),
\[
\boxed{
Q
\overset d=
q_0
+
\sum_i
\lambda_i
\chi^2_1(\delta_i),
}
\]
where
\[
\delta_i=
\left(\frac{c_i}{\lambda_i}\right)^2
\]
and
\[
q_0
=
\mu^\top M\mu
-
\sum_i\frac{c_i^2}{\lambda_i}
\ge0.
\]
The noncentral chi-square components are independent. To state the
zero-shift condition precisely, put $T=M^{1/2}S$, $b=M^{1/2}\mu$,
and let $P_T$ be the orthogonal projector onto $\operatorname{range}(T)$.
Then
\[
q_0=\|(I-P_T)b\|^2.
\]
Thus $q_0=0$ exactly when $M^{1/2}\mu\in\operatorname{range}(M^{1/2}S)$.
In particular, $\Sigma\succ0$ is sufficient. For singular covariance the
shift is the squared deterministic component perpendicular to the random
range in these weighted coordinates.
\end{theorem}

\begin{proof}
Expand
\[
(\mu+SZ)^\top M(\mu+SZ)
=
Z^\top SMSZ
+
2\mu^\top MSZ
+
\mu^\top M\mu
\]
and rotate by \(V\).  Completing the square gives the displayed
noncentral components.  For a null vector of \(SMS\),
\[
M^{1/2}Sv=0,
\]
hence
\[
v^\top SM\mu
=
(M^{1/2}Sv)^\top M^{1/2}\mu
=
0,
\]
so there is no unresolved stochastic linear term in a zero-eigenvalue
direction. For the positive eigenvectors $v_i$ of $T^\top T=SMS$,
the vectors $Tv_i/\sqrt{\lambda_i}$ form an orthonormal basis of
$\operatorname{range}(T)$. Their inner products with $b$ are
$c_i/\sqrt{\lambda_i}$. Hence the subtracted sum equals $\|P_Tb\|^2$,
leaving $q_0=\|(I-P_T)b\|^2\ge0$. Orthogonal rotation preserves
independence of the normal coordinates, so the squared components are
independent as stated.
\end{proof}

The numerical interface implements this reduction directly.  In an isotropic
control,
\[
\Sigma=I,\qquad M=aI,\qquad a>0,
\]
the law collapses to the exact scaled noncentral chi-square distribution
\[
Q/a\sim\chi^2_d(\|\mu\|^2).
\]

The numerical controls for this reduction appear in
Chapter~\ref{ch:evidence}. They test a floating analytic-ledger engine and
are not assigned the directed-enclosure guarantee of the central positive
backend.

\section{Decomposition axes}

The solver separates four logically different representations:
\[
\boxed{
\text{storage blocks}
\circ
\text{structural phases}
\circ
\text{Krylov/range compression}
\circ
\text{spectral tail packets}.
}
\]
Storage blocks answer where data reside.  Structural phases answer how the
operator propagates.  Krylov/range compression answers which subspace has
actually been resolved.  Spectral packets answer how the final transform is
evaluated.  These decompositions should not be conflated.




\section{A factor, a spectrum, and an exact law}\label{sec:worked-law}
We will follow one small example through the remaining parts of the book. Its entries are exact binary fractions, so the same matrix can be used for symbolic algebra and for a binary64 input without introducing an input-rounding ambiguity. Set
\[
A=\frac12
\begin{pmatrix}
 1& 1& 1& 1& 1& 1\\
 1&-1& 1&-1& 1&-1\\
 1& 1&-1&-1&-1&-1\\
 1&-1&-1& 1&-1& 1
\end{pmatrix}.
\]
Multiplication gives
\[
K=AA^\top=\frac12
\begin{pmatrix}
3&0&-1&0\\0&3&0&-1\\-1&0&3&0\\0&-1&0&3
\end{pmatrix}.
\]
The orthogonal matrix
\[
V=\frac12
\begin{pmatrix}
1&1&1&1\\1&-1&1&-1\\1&1&-1&-1\\1&-1&-1&1
\end{pmatrix}
\]
satisfies $V^\top K V=\operatorname{diag}(1,1,2,2)$. For $Z\sim N(0,I_4)$,
\[
Q=Z^\top KZ\ \overset d=\ Y+2W,
\qquad Y,W\text{ independent with law }\chi^2_2.
\]
The spectrum is therefore $(1,1,2,2)$ and $\mathbb E Q=6$. The six-dimensional feature Gram $A^\top A$ has these same four nonzero eigenvalues and two additional zeros.

Here the probability law also has a closed form. The densities of $Y$ and $2W$ are $\tfrac12e^{-x/2}$ and $\tfrac14e^{-x/4}$ on the nonnegative half-line. Their convolution is
\[
f_Q(x)=\int_0^x\frac12e^{-y/2}\frac14e^{-(x-y)/4}\,dy
       =\frac12\bigl(e^{-x/4}-e^{-x/2}\bigr).
\]
Integrating from $q$ to infinity gives
\begin{equation}
P(Q\ge q)=2e^{-q/4}-e^{-q/2},\qquad q\ge0.
\label{eq:worked-exact-tail}
\end{equation}
At $q=40$ the probability is approximately $9.079779837\times10^{-5}$. This formula will be used as an independent control, not as a shortcut inside the COS calculation. The next part begins with the factor $A$ and asks which operations may change its storage or basis without changing this law.
